\documentclass[a4paper,12pt]{article}
\usepackage[utf8]{inputenc}
\usepackage[T2A]{fontenc}
\usepackage[russian]{babel}
\usepackage{amsmath, amssymb, amsthm}
\usepackage{fullpage}
\usepackage{enumitem}

\newtheorem{theorem}{Теорема}[section]
\newtheorem{lemma}[theorem]{Лемма}
\title{}
\author{}
\date{}
\begin{document}
	
	\section*{Основная часть}
	

	
	\subsubsection*{Задача 1: Выпуклость функции $f(x) = x_1^2 x_2^2$}
	
	Рассмотрим функцию $f: \mathbb{R}^2 \rightarrow \mathbb{R}$, заданную как $f(x) = f(x_1, x_2) = x_1^2 x_2^2$. Для определения выпуклости функции, мы можем исследовать ее гессиан.
	
	Сначала найдем первые частные производные:
	\begin{align*}
		\frac{\partial f}{\partial x_1} &= 2x_1 x_2^2 \\
		\frac{\partial f}{\partial x_2} &= 2x_1^2 x_2
	\end{align*}
	
	Теперь найдем вторые частные производные:
	\begin{align*}
		\frac{\partial^2 f}{\partial x_1^2} &= \frac{\partial}{\partial x_1} (2x_1 x_2^2) = 2x_2^2 \\
		\frac{\partial^2 f}{\partial x_2^2} &= \frac{\partial}{\partial x_2} (2x_1^2 x_2) = 2x_1^2 \\
		\frac{\partial^2 f}{\partial x_1 \partial x_2} &= \frac{\partial}{\partial x_1} (2x_1^2 x_2) = 4x_1 x_2 \\
		\frac{\partial^2 f}{\partial x_2 \partial x_1} &= \frac{\partial}{\partial x_2} (2x_1 x_2^2) = 4x_1 x_2
	\end{align*}
	
	Гессиан функции $f$ имеет вид:
	$$
	H_f(x) = \begin{pmatrix}
		\frac{\partial^2 f}{\partial x_1^2} & \frac{\partial^2 f}{\partial x_1 \partial x_2} \\
		\frac{\partial^2 f}{\partial x_2 \partial x_1} & \frac{\partial^2 f}{\partial x_2^2}
	\end{pmatrix} =
	\begin{pmatrix}
		2x_2^2 & 4x_1 x_2 \\
		4x_1 x_2 & 2x_1^2
	\end{pmatrix}
	$$
	
	Для того чтобы функция $f$ была выпуклой, ее гессиан должен быть положительно полуопределенным для всех $x \in \mathbb{R}^2$. Для матрицы $2 \times 2$, это означает, что главные миноры должны быть неотрицательными.
	
	Первый главный минор:
	$$
	\Delta_1 = 2x_2^2 \ge 0 \quad \forall x_2 \in \mathbb{R}
	$$
	
	Второй главный минор (определитель гессиана):
	$$
	\Delta_2 = \det(H_f(x)) = (2x_2^2)(2x_1^2) - (4x_1 x_2)(4x_1 x_2) = 4x_1^2 x_2^2 - 16x_1^2 x_2^2 = -12x_1^2 x_2^2
	$$
	
	Для положительной полуопределенности гессиана, необходимо, чтобы $\Delta_2 \ge 0$. Однако, $-12x_1^2 x_2^2 \ge 0$ только в том случае, если $x_1 = 0$ или $x_2 = 0$. Например, если $x_1 = 1$ и $x_2 = 1$, то $\Delta_2 = -12 < 0$.
	
	Так как определитель гессиана может быть отрицательным, функция $f(x) = x_1^2 x_2^2$ не является выпуклой на $\mathbb{R}^2$.
	
	\subsubsection*{Задача 2:}
	
	Решение. Для каждой координаты функция $x \mapsto x^4$ имеет вторую производную:
	$$
	\frac{d^2}{dx^2} x^4 = 12x^2.
	$$
	Так как $12x^2 \ge 0$ для всех $x \in \mathbb{R}$, функция $x^4$ является выпуклой. Следовательно, сумма таких функций по разным координатам также выпукла.
	
	Для $\mu$-сильной выпуклости требуется, чтобы для всех $x \in \mathbb{R}^d$ гессиан $H_f(x)$ удовлетворял неравенству $H_f(x) \ge \mu I$.
	
	При вычислении гессиана получаем диагональную матрицу с элементами $12x_i^2$. Заметим, что при $x = 0$ все диагональные элементы равны нулю, то есть $H_f(0) = 0$. Таким образом, не существует $\mu > 0$, для которого $H_f(x) \ge \mu I$ для всех $x$.
	
	Ответ: функция $f(x) = \sum_{i=1}^{d}x_i^4$ выпукла, но не является $\mu$-сильно выпуклой ни для какого $\mu > 0$.
	
	\subsubsection*{Задача 3:}
	
	Решение. Из классических результатов выпуклого анализа и теории спектральных функций известно следующее:
	
	1) Функция, сопоставляющая симметричной матрице её наибольшее собственное значение, является выпуклой.
	$$
	f(X) = \lambda_{max}(X) \quad \text{выпуклая}
	$$
	
	2) Функция, сопоставляющая симметричной матрице её наименьшее собственное значение, является вогнутой.
	$$
	f(X) = \lambda_{min}(X) \quad \text{вогнутая}
	$$
	
	\subsubsection*{Задача 4:}
	
	Условие. Пусть $f: \text{dom } f \rightarrow \mathbb{R}$ - функция с областью определения $\text{dom } f \subseteq \mathbb{R}^d$. Докажите, что $f$ выпукла тогда и только тогда, когда её сужение на любую прямую является выпуклым. Формально: для любых $x_0 \in \text{dom } f$ и $u \in \mathbb{R}^d$ функция $g(t) = f(x_0 + tu)$, $t \in \{t \in \mathbb{R} : x_0 + tu \in \text{dom } f\}$, является выпуклой.
	
	Доказательство.
	
	(Необходимость). Пусть $f$ выпукла, т.е. для любых $x, y \in \text{dom } f$ и $\lambda \in [0, 1]$ выполнено:
	$$
	f(\lambda x + (1 - \lambda) y) \le \lambda f(x) + (1 - \lambda) f(y).
	$$
	Выберем произвольное $x_0 \in \text{dom } f$ и $u \in \mathbb{R}^d$. Для любых $t_1, t_2$ из области определения функции $g$ и для любого $\lambda \in [0, 1]$ рассмотрим точки $x_0 + t_1 u$ и $x_0 + t_2 u$. Тогда
	\begin{align*}
		g(\lambda t_1 + (1 - \lambda) t_2) &= f(x_0 + (\lambda t_1 + (1 - \lambda) t_2) u) \\
		&= f(\lambda(x_0 + t_1 u) + (1 - \lambda)(x_0 + t_2 u)) \\
		&\le \lambda f(x_0 + t_1 u) + (1 - \lambda) f(x_0 + t_2 u) \\
		&= \lambda g(t_1) + (1 - \lambda) g(t_2).
	\end{align*}
	Таким образом, $g$ выпукла.
	
	(Достаточность). Пусть для любого $x_0 \in \text{dom } f$ и любого направления $u \in \mathbb{R}^d$ функция $g(t) = f(x_0 + tu)$ выпукла на своем множестве определения. Возьмём произвольные $x, y \in \text{dom } f$ и положим $u = y - x$; тогда можно записать $x = x$ и $y = x + 1 \cdot u$. Определим функцию $g(t) = f(x + tu)$, которая по предположению выпукла для таких $t$, что $x + tu \in \text{dom } f$. При $t = 0$ и $t = 1$ имеем $g(0) = f(x)$ и $g(1) = f(y)$, а выпуклость $g$ даёт
	$$
	f(\lambda x + (1 - \lambda) y) = g(1 - \lambda) \le (1 - \lambda) g(0) + \lambda g(1) = (1 - \lambda) f(x) + \lambda f(y)
	$$
	для всех $\lambda \in [0, 1]$. Это и есть определение выпуклости функции $f$.
	
	Таким образом, $f$ выпукла тогда и только тогда, когда её сужение на любую прямую выпукло.
	
	\subsection*{Дополнительная часть}
	
\subsection*{Задача 1}
Пусть функция $\varphi(x) = -\ln x$, она выпукла на $(0,\infty)$. По неравенству Йенсена для весов $p_i\ge0$, $\sum_i p_i=1$ имеем
\[
-\ln\Bigl(\sum_{i=1}^d p_i\frac{q_i}{p_i}\Bigr)\le\sum_{i=1}^d p_i(-\ln\tfrac{q_i}{p_i}).
\]
Так как $\sum_i q_i=1$, левая часть равна $-\ln1=0$, отсюда
\[
0\le\sum_{i=1}^d p_i\ln\frac{p_i}{q_i}.
\]

\subsection*{Задача 2}
Пусть \(g\colon\mathbb R_{+}\to\mathbb R_{+}\) — выпуклая функция, причём \(g(0)=0\). Определим
\[
f(x)=\frac1x\int_0^x g(t)\,\mathrm{d}t,\qquad x>0.
\]
Сделаем замену переменной \(t=ux\), где \(u\in[0,1]\). Тогда
\[
f(x)
=\frac1x\int_0^x g(t)\,\mathrm{d}t
=\int_0^1 g(ux)\,\mathrm{d}u.
\]
Для каждого фиксированного \(u\in[0,1]\) отображение 
\[
x\;\longmapsto\;g(ux)
\]
является выпуклым, поскольку это композиция выпуклой функции \(g\) с аффинным отображением \(x\mapsto ux\). Интеграл по \(u\) от выпуклых функций сохраняет выпуклость. Следовательно, \(f(x)\) выпукла на \(\mathbb R_{++}\).

\subsection*{Задача 3}
Пусть $f(X)=\mathrm{Tr}(X^{-1})$ для $X\in S_{++}^d$. Дифференциал функции:
\[Df(X)[H]=-\mathrm{Tr}(X^{-1}HX^{-1}),\quad H=H^T.\]
Второй дифференциал в направлении $H$:
\[D^2f(X)[H,H]=2\,\mathrm{Tr}(X^{-1}HX^{-1}HX^{-1})\ge0,\]
так как матрица $X^{-1}HX^{-1}HX^{-1}$ положительно полуопределённая и её след неотрицателен. Следовательно $f$ выпукла на $S_{++}^d$.
	Ответ: функция $f(X) = Tr(X^{-1})$ выпукла на $\mathbb{S}_{++}^{d}$.
	
	\subsubsection*{Задача 4}
	
	Воспользовавшись неравенством Йенсена для выпуклой на $\mathbb{R}_{++}$ функции $f(x) = -\ln x$, докажите неравенство Гёльдера:
	для $p > 1$, $\frac{1}{p} + \frac{1}{q} = 1$.
	$$
	\sum_{i=1}^{d} x_i y_i \le \left(\sum_{i=1}^{d} |x_i|^p\right)^{1/p} \left(\sum_{i=1}^{d} |y_i|^q\right)^{1/q}.
	$$
	
	Решение.
	
	Шаг 1. Для набора положительных чисел $\{a_i\}$ и весов $\{\alpha_i\}$, где $\alpha_i \ge 0$, $\sum_{i=1}^{d} \alpha_i = 1$, неравенство Йенсена для выпуклой функции $-\ln x$ гласит:
	$$
	-\ln\left(\sum_{i=1}^{d} \alpha_i a_i\right) \le \sum_{i=1}^{d} \alpha_i (-\ln a_i).
	$$
	Эквивалентно,
	$$
	\ln\left(\sum_{i=1}^{d} \alpha_i a_i\right) \ge \sum_{i=1}^{d} \alpha_i \ln a_i.
	$$
	
	Шаг 2. Выберем
	$$
	\alpha_i = \frac{|x_i|^p}{\sum_{j=1}^{d} |x_j|^p}, \quad a_i = \frac{|y_i|}{|x_i|^{p/q}},
	$$
	при условии, что все $x_i \ne 0$ (остальные случаи рассматриваются аналогично, а предел можно получить переходом к пределу).
	
	Тогда неравенство Йенсена приводит к соотношению:
	$$
	\ln\left(\sum_{i=1}^{d} \frac{|x_i|^p}{\sum_{j=1}^{d} |x_j|^p} \frac{|y_i|}{|x_i|^{p/q}}\right) \ge \sum_{i=1}^{d} \frac{|x_i|^p}{\sum_{j=1}^{d} |x_j|^p} \ln\left(\frac{|y_i|}{|x_i|^{p/q}}\right).
	$$
	Учитывая, что $\frac{p}{q} = p(1 - \frac{1}{p}) = p - 1$, имеем $|x_i|^{p/q} = |x_i|^{p-1}$.
	$$
	\ln\left(\frac{\sum_{i=1}^{d} |x_i|^{p - (p-1)} |y_i|}{\sum_{j=1}^{d} |x_j|^p}\right) \ge \frac{1}{\sum_{j=1}^{d} |x_j|^p} \sum_{i=1}^{d} |x_i|^p (\ln |y_i| - (p-1) \ln |x_i|).
	$$
	$$
	\ln\left(\frac{\sum_{i=1}^{d} |x_i| |y_i|}{\sum_{j=1}^{d} |x_j|^p}\right) \ge \frac{\sum_{i=1}^{d} |x_i|^p \ln |y_i|}{\sum_{j=1}^{d} |x_j|^p} - (p-1) \frac{\sum_{i=1}^{d} |x_i|^p \ln |x_i|}{\sum_{j=1}^{d} |x_j|^p}.
	$$
	Применяя экспоненту к обеим частям и после преобразований (учитывая соотношение $\frac{1}{p} + \frac{1}{q} = 1$), получаем неравенство Гёльдера:
	$$
	\sum_{i=1}^{d} |x_i y_i| \le \left(\sum_{i=1}^{d} |x_i|^p\right)^{1/p} \left(\sum_{i=1}^{d} |y_i|^q\right)^{1/q}.
	$$
	Для действительных чисел $x_i, y_i$, имеем $|x_i y_i| = |x_i| |y_i|$, и $\sum_{i=1}^{d} x_i y_i \le \sum_{i=1}^{d} |x_i y_i|$.
	
\end{document}